The electronic structure theory method discussed in \cref{chapter:energy_models} have been the workhorse for modern materials discoveries, enabling a wide range of materials properties to be computed from first-principles [\cref{chapter:physical_properties}], that ultimately led to the establishment of various large-scale computational materials databases [\cref{chapter:high-throughput}] that revolutionised how we conduct theoretical materials research in the past decade. However, the power scaling laws of these theoretical methods (\cref{tab:scaling_low}) have fundamentally limited the system sizes that can be practically simulated, typically to a few hundred atoms for DFT and only tens of atoms for more accurate many-body perturbation theory methods such as GW. This has imposed a limit on the range of chemical and configurational spaces that can be explored. It also constrained the thermodynamic conditions $(T, P, V)$ under which we can perform statistical samplings to obtain physical observables close to experimental conditions. 

\begin{table}[htb]
    \centering
    \caption{Scaling behaviours of the theoretical methods for materials modelling discussed in this review. The system sizes are measured in terms of the number of atoms $N$ in the systems being simulated.}
    \label{tab:scaling_low}
    \begin{tabular}{llll}
    \hline\hline 
    \textbf{Method} & \textbf{Formal Scaling} & \textbf{Limiting Step} & \textbf{Practical System Sizes} \\
    \hline
    Classical forcefields & $\mathcal{O}(N)$ & for short-range interaction defined by cut-off & up to 10$^6$ \\
                         & $\mathcal{O}(N\log N)$ & for long-range Ewald summation & \\
    Pure DFT  & $\mathcal{O}(N^3)$ & matrix diagonalisation for KS eigenvalues & $\sim 10^{2}-10^{3}$\\
    Hybrid DFT  & $\mathcal{O}(N^4)$ (Na\"{i}ve) & exact exchange integral evaluation & $\sim 10^{2}$ \\
                & $\mathcal{O}(N^3)$ (Practical) & screened long-range exchange & \\
    $G_{0}W_{0}$ & $\mathcal{O}(N^4)$ & computation of the polarisation $\chi_0(\bm{r},\bm{r}^\prime)$ & $\sim 10-10^{2}$ \\
                                   && Inversion of the dielectric matrix &\\
    self-consistent $GW$ & $\mathcal{O}(N^4\sim N^5)$ & repeated calculations of $G$ and $W$ & $<50$ \\
    \hline\hline
    \end{tabular}
\end{table}

Even with the rapid advancement in HPC, particularly in GPU (Graphic Processing Unit)-accelerated electronic structure theories, there remain many challenging areas in which calculations using currently available first-principle methods are far from routine practices. These include, but are not limited to: (a) the determination of accurate thermodynamic properties for alloys, defects and disordered systems with the explicit account of the lattice vibrational contributions,\cite{van2002effect} (b) many-body physics problems such as electron-phonon interactions (\cref{sec:eph_coupling})\cite{giustino2017electron} in polar semiconductors and high-temperature superconductors, (c) free energy sampling for solvated surfaces and heterogeneous catalysis under explicit solvent at finite temperatures,\cite{gono2019effect,sakong2015density,gross2022ab} (d) vdW heterostructures of 2D materials with incommensurate lattice parameters  and Moir{\'e} pattern superlattices, which often require the use of large simulation supercells,\cite{tang2020moire} and finally, (f)  determining the properties, such as the lattice thermal conductivities\cite{allen1993thermal} and catalytic activities\cite{sun2021high} of high-entropy materials, particularly taking into account the configurational entropic contributions through ensemble averaging [\cref{eq:general_sampling_2}]. 

From a practicality perspective, although empirical atom–atom forcefields (\cref{sect:empirical_forcefield}) can efficiently simulate systems with millions of atoms, their limited transferability and accuracy often hinder predictive materials modelling across diverse chemical and structural spaces. In particular, fitting such forcefields to DFT data is generally a non-trivial process. This is because standard DFT does not decompose the total energy into physically interpretable contributions from different physical terms (such as electrostatic, dispersions and exchange-repulsions) as well as their corresponding atom-atom contributions. As a result, parameters from different forcefield terms become strongly coupled. The rigid functional forms, combined with limited training data, also yield narrow and stiff optimisation landscapes, which usually result in poor numerical stability during the parameter fitting procedures. These turn forcefield fitting into a craft rather than a robust, reproducible procedure. One remedy, for example, when fitting intermolecular interaction potentials, is to use SAPT to decompose the interaction energy into individual physical components prior to fitting [see \citet{misquitta2005intermolecular,misquitta2016ab}], although the overall workflow is still far from trivial. In this regard, the rise of big-data in materials science and the progressive wide adoption of the ML/AI techniques, which are established on a probabilistic framework that naturally encourages the introduction of noise during the learning process,  have provided new opportunities to address these challenges.

From a fundamental perspective, uncovering and understanding the complex chemistry-structure-property relationships in materials to guide their future discoveries and development is an ongoing and prominent challenge in computational materials science, particularly with the continuously growing amount of data available from energy landscape sampling (\cref{chapter:sampling_methods}) and modern high-throughput computations (\cref{chapter:high-throughput}). This is because the  chemistry-structure-property relationship is often a high-dimensional function that is difficult for human to comprehend: the spatial arrangement of $N$ atoms in a molecule/materials leads to $3N$ degrees-of-freedoms, and from the viewpoint of DFT, the energy and physical properties of an atomistic system depends non-linearly on its ground-state electron density $\rho(\bm{r})$, which is also a geometric object with complicated shape in the 3D space. Consequently, the mapping from the atomic coordinates $\{\bm{R}\}$ to the electron density $\rho(\bm{r})$ and further to the total energy $E[\rho(\bm{r})]$ and other physical properties $\mathcal{P}[\rho(\bm{r})]$ are highly non-linear and high-dimensional functions that are difficult to be expressed in closed-form equations. Traditionally, this problem was overcome by the application of heuristic rules that were established based on careful human inspections on available materials data. Examples of these heuristic rules include (a) tolerance factors for rationalising the phase stabilities of ABX$_3$-type of perovskites,\cite{Goldschmidt1926_Krystallochemie,Shannon1976_IonicRadii,Bartel2019_NewToleranceFactor} (b) Pauling's rule for rationalising the close-packing and stabilities of inorganic crystals, (c) hydrogen bonding patterns and the stabilities of molecular crystals,\cite{gavezzotti2007molecular}   (d) $d$-band theory for rationalising the performance of transition-metal-based catalysts,\cite{Hammer1995_Nature_dband} and (e) suppression of lattice thermal conductivities due to mass and structural disorders in random alloys.\cite{Abeles1963_AlloyThermalConductivity} These heuristic rules are either quantitative or semi-qualitative, simple and intuitive for post-hoc analysis, but they are often limited in their accuracy and transferability, particularly when applied to predictive investigations on material systems that are outside the original domains in which they were devised. 

In this regard, the adaptation of ML/AI techniques in materials science during the past decade has brought a paradigm shift that enriched our capabilities in performing computational modelling on a more diverse range of materials. This lies critically on the capabilities of ML/AI models to learn and generate workable mathematical relationships among the chemistry, structure, energetic stabilities and physical properties for materials curated from structure predictions and high-throughput computations with little or no human bias. Direct utilisation of these mathematical relationships learned by the models provides orders-of-magnitude acceleration in computational materials modelling while retaining the accuracy of high-level electronic structure theories, thus circumventing the limitations imposed by the scaling laws. Augmented with advanced analysis and visualisation techniques further enables us to explore these complex structure-property relationships in an intuitive and visually appealing way to enhance our fundamental understanding on the behaviours of materials at the atomistic level. Finally, recent advancements in generative AI have also opened new opportunities for inverse materials design by holistically harnessing the complex chemistry-structure-property relationships from a diverse range of materials, where the desired properties are specified as inputs to generate candidate materials structures that can meet these property requirements. These advancements are inevitably fueled by the wide range of computational materials databases available, which provide consistent baselines for model training and validations.

Before diving into the details, we shall highlight some general conceptual points regarding ML and AI, where some of them are more specifically applicable to materials science.

Firstly, it is worth noting that ML and AI are often used interchangeably in the literature, and the techniques developed for ML and AI often complement each other in practice. In this review, particularly with respect to the research in materials science, we distinguish ML and AI based on the following general characteristics. For ML, we consider it as an approach that generally relies on some hand-crafted mathematical descriptors constructed from physical insights for representing materials' structures\cite{musil2021physics} and properties. Consequently, ML algorithms typically work on small to medium datasets to generate simple (often of linear characters) models with strong physical interpretability. On the other hand, AI or deep learning (DL, which we shall be using interchangeably from now on) from raw or minimally processed data to generate latent features (analogous to descriptors from ML) that link the input with the output.  DL models, through the implementations of NNs that were designed to mimic the way information is passed and processed through the interconnected neurons in the human brains,  usually possess complex model architectures with thousands or even millions of model parameters that need to be trained on a large amount of data. Whilst the increased complexity in the DL models often leads to better generalisability and higher accuracy, their physical interpretability can be significantly compromised, and very often the trained model can only be considered as a valid \emph{ansatz} within the domain of the training data and the chosen hyperparameters for the specific model architecture.   

Secondly, parallel to ML and AI are the concepts of supervised and unsupervised learning. Built upon some material representation $\mathcal{X}$ (which can be hand-crafted for ML or learned by the DL models), \textbf{\textit{supervised learning}} tries to learn a parameterised function $y=f_{\boldsymbol{\theta}}(\mathcal{X})$ that maps the input $\mathcal{X}$ to a known output $y$ (e.g., material properties) based on a set of training data $\{(\mathcal{X}_i, y_i)\}$, where $\boldsymbol{\theta}$ are the model parameters to be optimised during the training process. On the other hand, \textbf{\textit{unsupervised learning}} aims to discover the relationships amongst the materials (especially their structures) based on the similarities between materials, which are typically measured as some distance functions based on $\mathcal{X}$. It is important to emphasise that supervised and unsupervised learning do not necessarily work in isolation, but can be  applied in conjunction with each other, for example, in refining the material representation $\mathcal{X}$ to deliver the best learning outcome for both scenarios.\cite{helfrecht2020structure} 

Finally, different from the tasks such as image recognition and language modelling in AI, there are many existing physical and chemical rules which dictate the behaviours of molecules and materials at the atomistic scales. Rather than letting the models to discover these physical rules themselves, which could exponentially increase the amount of training data required, modern ML/AI seeks to incorporate these rules into the model as an \emph{inductive bias} or \emph{physical constraints} to reduce the training cost, improve the model's stability and generalisability. However, this sacrifices the model complexities, and implementing these architectures often requires deep and specialised mathematical knowledge in areas such as group and representation theories, as well as tensor algebra, which are not in the knowledge bases of common practitioners of computational materials.  This is usually coined as \emph{physically informed} ML/AI for materials. Examples of these physical constraints include (a) translational invariance and rotational equivariance of the physical observables with respect to the atomic positions,\cite{Bartok2013_RepresentingChemEnv,grisafi2018symmetry,duval2023hitchhiker} (b) conservation of energy and momentum in molecular dynamics simulations,\cite{chmiela2017machine} (c) crystal point-group and space-group symmetry for generating stable crystal structures,\cite{jiao2024space,zhu2024wycryst} (d) long-ranged atom-atom interactions in chemical structures.\cite{grisafi2019incorporating,frank2022so3krates} These will be discussed in more detail later. 

This part of the review is intended to provide an account for the core and generic concepts behind modern ML/AI that are highly relevant to the research in materials science to facilitate the development of a fundamental understanding of the field. 
When possible, we will try to establish some insightful connections among ML/AI methodologies with traditional atomistic modellings (e.g. forcefield methods) and energy landscape samplings. More specifically, this part of the review is organised as follows. 

\cref{sec:prob_theory_fundamentals} provides a brief overview of the fundamentals of probability theory that underpins modern ML/AI techniques.\cite{murphy2012machine,goodfellow2016deep} Such a viewpoint shares many conceptual similarities to the ensemble-averaging view of the structure-property landscapes for materials [\cref{eq:general_sampling_1} and \cref{eq:general_sampling_2}]. 

\cref{sec:ml_fundamentals} introduces the core concepts behind,  and techniques commonly applied for ML, including those for supervised and unsupervised learning. 

\cref{sec:feature_eng} discusses the construction of hand-crafted descriptors for representing materials' structures. We devote this section to discussing this important aspect in ML for materials not only because it plays a pivotal role behind the success of ML, but it also provides a strong foundation for understanding the incorporation of physical constraints in DL models to enhance the model stabilities and generalisabilities. 

\cref{sec:dl_fundamentals} will be focusing on the key concepts in modern DL architectures and model training, in particular, how the model architectures mimic our conventional understanding of atom-atom interactions in molecular and materials systems. We shall also discuss the importance of latent space in DL to enhance the model interpretability for facilitating the development of fundamental understanding. In this part of discussions around DL, we aim to avoid elaborate discussion on the architectural and implementation choices for specific models, as some of most essential mathematical frameworks behind the model designs are rather consistent amongst different models being designed. 

\cref{sec:generative_AI} discusses the recent advancement in generative AI for materials design with special attention being paid to the incorporation of physical constraints in materials generation. Here, we focus specifically on the development of variational autoencoders, diffusion and normalising-flow generative frameworks. Unlike many other sections in this monograph, which aims to provide a general overview on the theoretical methodologies, we have deliberately chosen  to provide more detailed account on specific models that have been developed under each framework. This is because we feel that the model design space, encompassing interdisciplinary knowledges in structural representation, symmetry constraints, statistical sampling, stochastic dynamics, differential geometries, group theory and beyond, is enormously large, that each one of the as-developed model is a unique and intriguing combination of these aspects, with much of the development being driven by computer scientists, rather than traditional practitioners of atomistic modellings. Given the recency of generative AI models, we thus believe a more centralised detailed account on the architectural designs behind these model will facilitate the atomistic modelling community to gain a better understanding of this rapidly developing field.

%\cref{sec:application_AI} provides an account for the wide range of applications of ML/AI in materials science, including structure predictions, property predictions, interatomic potentials, and so on. 
%Finally, we highlight some immediate challenges that the community is facing with the application of ML/AI for materials studies in \cref{sec:challenge_AI}.  

